\section{Conclusion}
%"novel consensus layer that finalizes transactions with the optimal one-round-trip latency" sounds a bit fluffy, let's make this more precise for the DISC submission, or just remove this Discussion section.
In this work we present \podtm, a novel consensus layer that finalizes transactions with the optimal one-round-trip latency by eliminating communication among replicas. Instead, clients read the system state by performing lightweight computation on logs retrieved from the replicas. As no validator has a particular role in \podtm (as compared to leaders, block proposers, miners, etc. in similar protocols), \podtm achieves censorship resistance by default, without any extra mechanisms or additional cost. Furthermore, validator misbehavior, such as voting in incompatible ways or censoring confirmed transactions, is accountable.

As an application, we present an efficient and censorship-resistant mechanism for single-shot first-price and second-price open auctions, which leverages \podtm as a bulletin board.
We show how the accountability, offered by \podtm, is also inherited by applications built on it -- the auctioneer cannot censor confirmed bids without being detected.
In the full version~\cite{DBLP:journals/corr/abs-2501-14931} of this paper we discuss further applications, including payments in the style of Fastpay~\cite{DBLP:conf/aft/BaudetDS20}.

We conjecture that single-shot sealed bid auction protocols, such as those of~\cite{trustee,seal,ACNS:DavGenPou22,CCS:TAFWBM23,chitra_et_al:LIPIcs.AFT.2024.19,cryptoeprint:2024/1011}, can also be instantiated on top of a \bidsettm protocol. Intuitively, this holds because such protocols first agree on a set of sealed bids and then determine the winner. A formal analysis of sealed-bid auction protocols based on \bidsettm is left as future work.
The architecture of \podtm also motivates the design of light clients, which would not have to connect to all replicas or to download all transactions. This is achieved through cryptographic primitives such as Merkle mountain ranges and segment trees. We leave the formal description as future work.

Finally, we remark that \podtm differs from standard notions of consensus because it does not offer an agreement property, neither to validators nor to clients. A client reading the \podtm obtains a past-perfect round \rperf, and it is guaranteed to have received all transactions that obtained a confirmed round \rconf such that $\rconf \leq \rperf$, or that may obtain such an $\rconf$ in the future, even though the transaction presently appears to the client as unconfirmed. However, the client cannot tell which unconfirmed transactions will become confirmed. Moreover, a transaction might appear confirmed to one client and unconfirmed to another (in this case, this is a transaction written by a malicious client).